% Part: normal-modal-logic
% Chapter: tableaux
% Section: introduction

\documentclass[../../../include/open-logic-section]{subfiles}

\begin{document}

\olfileid{nml}{tab}{int}

\olsection{مقدمة}

ال!!^{tableau}s أشجار مخصوصة تتفرع إلى أسفل، وعُقَدها !!{signed
  formula}s. والمراد بذلك أزواج تجمع علامة صدق ($\True$ أو
$\False$) و!!a{sentence}، على أحد الوجهين:
\[
\sFmla{\True}{!A} \text{ أو } \sFmla{\False}{!A}.
\]
نبتدئ !!^a{tableau} بجملة من \emph{الافتراضات}، ثم لا نزيد
!!{signed formula} إلا بقاعدة من قواعد الاستدلال. فمن القواعد ما يلحق
!!{signed formula}s واحدة أو أكثر بطرف الشجرة، ومنها ما ينشئ طرفين
جديدين، فتتشعب الشجرة شعبتين. وتكون ال!!{signed formula}s الناتجة ذات
!!{formula} أقل تعقيدًا مما في ال!!{signed formula} التي أُجريت عليها
القاعدة. فإذا اجتمع في شعبة~$\sFmla{\True}{!A}$ و$\sFmla{\False}{!A}$
قلنا إنها \emph{مغلقة}. ومتى أُغلقت شعب !!a{tableau} كلها أُغلقت
ال!!{tableau} بأسرها. فال!!{tableau} المغلقة !!a{derivation} تثبت أن
مجموعة ال!!{signed formula}s التي افتُتحت بها ال!!{tableau} غير قابلة
للإرضاء. وبهذا نعرّف العلاقة~$\Proves$: يكون~${\OLId{greek-variable}{Gamma}} \Proves !A$
إذا وفقط إذا وجدت مجموعة منتهية
${\OLId{greek-variable}{Gamma}}_{\OLMathNumeral{0}}=\{!B_{\OLMathNumeral{1}},\dots,!B_{\OLId{latin-ordinary}{n}}\}\subseteq{\OLId{greek-variable}{Gamma}}$ لها !!{tableau} مغلقة
تبدأ بالافتراضات:
\[
\{\sFmla{\False}{!A}, \sFmla{\True}{!B_{\OLMathNumeral{1}}}, \dots, \sFmla{\True}{!B_{\OLId{latin-ordinary}{n}}}\}.
\]

وأما المنطق الجهوي فيحتاج إلى توسيع معنى !!{signed formula} وإلى
قواعد للمؤثرات~\iftag{prvBox}{$\Box$\iftag{prvDiamond}{ و
    $\Diamond$}{}}{$\Diamond$}. فلا تقتصر ال!!{formula}s، بعد العلامة
($\True$ أو~$\False$)، في ال!!{tableau}s الجهوية على ذلك، بل تلحقها
\emph{بادئات}~$\sigma$، وهي متتاليات غير خالية من الأعداد الصحيحة
الموجبة؛ أي~$\sigma \in (\PosInt)^*\setminus\{\emptyseq\}$.
ونكتبها من غير القوسين~$\tuple{\ }$، ونجعل بين ال!!{element}s المفردة
نقطًا~$.$ عوض الفواصل~$,$. فإذا كانت~$\sigma$ بادئة كان~$\sigma.{\OLId{latin-ordinary}{n}}$
اختصارًا لـ~$\sigma\concat\tuple{{\OLId{latin-ordinary}{n}}}$؛ فإن كان~$\sigma={\OLMathNumeral{1}}.{\OLMathNumeral{2}}.{\OLMathNumeral{1}}$ مثلًا،
كان~$\sigma.{\OLMathNumeral{3}}$ هو~${\OLMathNumeral{1}}.{\OLMathNumeral{2}}.{\OLMathNumeral{1}}.{\OLMathNumeral{3}}$. وعلى هذا تكون
\[
\sFmla{\True}{\Box !A \lif !A}[{\OLMathNumeral{1}}.{\OLMathNumeral{2}}]
\]
هي \emph{!!{signed formula} ذات بادئة} (أو اختصارًا \emph{!!{formula}
  ذات بادئة}).

ومأخذ هذا الاصطلاح أن البادئة اسم لعالم في نموذج يُرجى أن يرضي
ال!!{formula}s الواقعة على شعبة من !!a{tableau}. فإن سمّت~$\sigma$
عالمًا، سمّت~$\sigma.{\OLId{latin-ordinary}{n}}$ عالمًا ينفذ إليه العالم المسمى بـ~$\sigma$.

\end{document}
